\chapter{大自然は「究極の怠け者」である --- 最小作用の原理と大域的最適化}

\section*{本章の学習目標}
\begin{itemize}
\item 光の屈折に関する「フェルマーの原理」から、自然界に潜む「目的論的」な振る舞いのミステリーを知る。
\item ニュートン力学（ベクトル）から解析力学（エネルギーと仕事）へのパラダイムシフトを理解する。
\item 物理学における最も美しく深遠な法則「最小作用の原理」の概念を学ぶ。
\item デカルト座標系の限界と「一般化座標」による物理法則の解放について知る。
\item ハミルトニアンと「位相空間」の概念が、AIの高次元空間探索とどう結びつくかを理解する。
\item 局所的な一歩（勾配降下法）に頼る現在のAIが抱える弱点と、全体を見渡す大域的最適化（強化学習など）の関連性を考察する。
\end{itemize}

\section{フェルマーの原理と「光の意志」のミステリー}

第1章で見たニュートンの運動方程式（\(F=ma\)）は、「今、この瞬間に背中を押されたから、次の瞬間に前へ進む」という、極めて常識的な「原因と結果（因果律）」に基づくルールでした。

しかし、17世紀のヨーロッパでは、もう一つ、全く別の視点から自然界のルールを紐解こうとする奇妙な動きが始まっていました。その発端となったのが、フランスの数学者ピエール・ド・フェルマーが発見した\textbf{「光の最短時間の原理（フェルマーの原理）」}です。

\subsection{デカルトとの論争と「自然の経済学」}

水の中に差し込んだストローが曲がって見えるように、光は空気中から水中に入る境界で「屈折」します。17世紀前半には、この曲がる角度のルール（スネルの法則）自体はすでに発見されていました。しかし、「なぜ光はその特定の角度で曲がるのか」という理由は謎に包まれていました。

同時代の偉大な哲学者ルネ・デカルトは、光を「力学的な粒子」として扱い、この屈折を説明しようとしました。しかし彼の理論には、「光は空気中よりも、抵抗が大きいはずの水中やガラスの中の方が『速く』進む」という、直感に反する強引な仮定が含まれていました。

フェルマーは、このデカルトの不自然な説明に納得がいきませんでした。「抵抗の大きい水中では、光のスピードはむしろ『遅くなる』はずだ」。そう考えたフェルマーは、古代から続く\textbf{「自然は常に最短かつ最も容易な道を選ぶ（自然の経済学）」}という哲学的な信念を、光の軌道に当てはめてみたのです。

彼自身が独自に編み出していた「極大・極小を求める数学的手法（ニュートンより少し前に微積分の先駆けとなったテクニック）」を用いて計算した結果、フェルマーは1662年に驚くべき真理へと到達しました。

「光は出発点から到着点まで移動する際、単に距離が短い直線ではなく、無数にある経路の中から『最も所要時間が短くなるルート』をあらかじめ選んで進んでいる」という解釈を与えたのです。

ここで大切なのは、光が本当に頭の中で全経路を比較している、という意味ではないことです。フェルマーの原理は、観測される光の経路を後から眺めると「所要時間が最小になる条件」を満たしている、という数学的な記述です。しかしこの記述の仕方があまりにも見事だったため、物理学者たちは「自然はまるで全体を知っているかのように振る舞う」という深い謎に引き込まれていきました。

\subsection{ライフセーバーの例えと、忍び寄る「目的論」}

有名な「ライフセーバーの例え」で考えてみましょう。

砂浜にいるライフセーバーが、海で溺れている人を助けに行くとします。砂浜を走るスピードは速いですが、海の中を泳ぐスピードは遅くなります。このとき、溺れている人へ向かって「直線の最短距離」を進むと、スピードの遅い海の中を長く泳がなければならず、結果的に到着が遅れてしまいます。少し遠回りになっても、砂浜を走る距離を長めに取り、海を泳ぐ距離を短くする「絶妙に折れ曲がったルート」を通るのが、最も早く到着する正解（最短時間）です。

光は、空気中（速い）からガラスや水中（遅い）に入るとき、まさにこのライフセーバーと同じように、全体を通して一番時間がかからないようにルートを曲げて進みます。

しかし、ここで恐ろしい哲学的な疑問が生まれます。

「光は、なぜゴール（到着点）の位置をあらかじめ知っているのか？」

「光は、無数にあるルートの所要時間をすべて計算した上で、出発する前に進路を決めているとでも言うのか？」

まるで光が「早くゴールに着きたい」という意志を持っているかのような、この\textbf{「目的論的」な振る舞いは、一歩ずつ因果律で進むニュートンの機械論的な宇宙観とは真っ向から対立するように見える、巨大なミステリーでした。 そして実は、この光が見せた奇妙な振る舞いこそが、のちに物理学全体を根本から書き換え、現代のAIの最適化概念にまで通じることになる「最小作用の原理」の最初の影}だったのです。

\section{解析力学の誕生：神学から数学、そして「エネルギー」の支配へ}

フェルマーが光について発見したこのミステリーを、「物質の運動」や「宇宙全体」へと拡張しようとする試みは、18世紀に入り、劇的な進化を遂げます。ここから、力学のあり方を根本から覆す\textbf{「解析力学（Analytical Mechanics）」}という壮大な理論が誕生します。

\subsection{モーペルテュイの神学的直感と「最小作用の原理」}

発端となったのは、フランスの哲学者・数学者であるピエール・ルイ・モーペルテュイでした。彼は1744年、ニュートン力学を深く研究する中で、ある哲学的な信念に行き着きます。

「神が創造したこの宇宙は、無駄を嫌い、完璧に経済的にできているはずだ。ならば、光だけでなく、石の落下や惑星の運動に至るまで、大自然のあらゆる現象は、常に『最も少ない労力（作用）』で物事を行おうとするはずだ。」

モーペルテュイはこの神学的・哲学的な直感を\textbf{「最小作用の原理」}と名付けて提唱しました。しかし、彼が言う「作用（労力）」の数学的な定義は曖昧であり、この美しいアイデアを真の物理学へと昇華させるためには、さらなる天才たちの登場を待たねばなりませんでした。

\subsection{ニュートン力学の「扱いづらさ」とベクトルの呪縛}

なぜ、すでに完璧なニュートンの運動方程式（\(F=ma\)）があるのに、新しい力学を作る必要があったのでしょうか？ それは、ニュートン力学が\textbf{「力（ベクトル）」}をベースにしているため、現実の複雑な問題を解く際に極めて「扱いづらい」からでした。

ベクトルは「大きさと向き」を持ちます。例えば、グネグネと曲がったレール（ジェットコースター）の上を滑り落ちるボールの動きをニュートン力学で計算しようとすると、重力だけでなく、レールがボールを押し返す力（垂直抗力）を、斜面の角度に合わせていちいち矢印に分解し、上下左右の成分を一つ残らず正確に足し合わせなければなりません。関節が複数ある振り子（多重振り子）ともなれば、その矢印の計算は絶望的なほど複雑に絡み合ってしまいます。

\subsection{ラグランジュの飛躍：ベクトル（力）からスカラー（エネルギーと仕事）へ}

そこで、フランスの天才数学者ジョゼフ＝ルイ・ラグランジュらは、向きを考えるのが面倒な「力（ベクトル）」を一切捨てるという飛躍に出ました。代わりに彼らが目をつけたのが、\textbf{「エネルギー（Energy）」と「仕事（Work）」}という概念です。

日常会話でも「仕事」という言葉を使いますが、物理学における「仕事」には極めて厳密な定義があります。それは\textbf{「物体に力（\(F\)）を加えて、その力の方向に距離（\(x\)）だけ移動させること（\(W=Fx\)）」}です。

例えば、重いボウリングの球を重力に逆らって高い棚の上に持ち上げる行為は、物理学的に「仕事をした」ことになります。

そして、この「仕事」とコインの表裏の関係にあるのが「エネルギー」です。エネルギーとは一言で言えば\textbf{「仕事をする能力（ポテンシャル）」}のことです。

あなたが汗をかいてボウリングの球を棚に持ち上げた（仕事をした）とき、その労力は消えてなくなるわけではありません。球の中に「高さ」という形で見えない能力（位置エネルギー）として蓄えられます。そして、その球が棚から転がり落ちると、今度は猛スピード（運動エネルギー）を獲得し、床に置かれた空き箱を粉々に吹き飛ばす（他の物体に対して仕事をする）能力を発揮するのです。

力（ベクトル）には「上下左右の向き」がありますが、エネルギーや仕事には「向き」がありません。温度や貯金残高と同じように、ただ「どれくらいの量があるか」だけを表す単なる数値（スカラー量）です。向きを気にしなくてよいスカラー量を使えば、どんなに複雑に曲がりくねったジェットコースターのレールでも、「高いところから低いところへ移動した（位置エネルギーが減った）分だけ、スピードが上がる（運動エネルギーが増える）」という足し算と引き算だけで、計算が劇的にシンプルになります。

ラグランジュが注目したのは、物体が持つ2種類の基本的なエネルギーでした。

運動エネルギー（\(K\)）：物体が「動くスピード」によって持っているエネルギー。（例：猛スピードで飛んでくるボールが持つ破壊力）

位置エネルギー（\(U\)）：物体が「高い場所」にあることなどで持っている、潜在的なエネルギー。（例：高い棚の上にある球が、落ちれば大きな破壊力を生む能力）

ラグランジュは、この2つのエネルギーの「差」を表す、のちに\textbf{「ラグランジアン（\(L\)）」}と呼ばれる指標を作りました。

\[
L=K-U
\]

そして、物体がスタート地点からゴール地点まであるルートを通って移動した際、その道筋に沿ってこのラグランジアンをすべて足し合わせた（時間積分した）合計値を、モーペルテュイが夢見た\textbf{「作用（Action：\(S\)）」}として厳密に数学的に定義したのです。

\[
S=\int_{t_1}^{t_2}L\,dt
\]

\subsection{変分法と「オイラー＝ラグランジュ方程式」の誕生}

ラグランジュと彼の師匠であるレオンハルト・オイラーは、この「作用\(S\)が最小になるルート」を探し出すために、\textbf{「変分法（Calculus of Variations）」と呼ばれる全く新しい魔法のような数学を創り上げました。 考えうる無限のルートの中から、作用\(S\)が最も小さく、あるいは少し経路をずらしても値が変わらない特別なルートを見つけ出したとき、そのルート上の物体は必ず次のような美しい微分方程式に従って動いていることが数学的に証明されました。これを「オイラー＝ラグランジュ方程式」}と呼びます。

厳密には、物理でよく言う「最小作用」は、常に文字通りの最小値を意味するわけではありません。実際には「作用が停留する」、つまり経路をほんの少し変えても作用の値が1次の範囲では変化しない、という条件です。山の頂上でも谷底でも鞍点でも、足元の傾きがゼロなら「停留点」です。このニュアンスを押さえると、最小作用の原理は神秘的な標語ではなく、微分で表せるきわめて精密な数学条件として理解できます。

\[
\frac{d}{dt}\left(\frac{\partial L}{\partial v}\right)-\frac{\partial L}{\partial x}=0
\]

（※\(x\)は位置、\(v\)は速度、\(L\)はラグランジアン）

この数式は一見複雑に見えますが、直感的な意味は極めてシンプルです。

右側の\(\frac{\partial L}{\partial x}\)は「位置を少しずらしたときに、エネルギーの無駄（\(L\)）がどう変化するか」を見ています。一方、左側の\(\frac{\partial L}{\partial v}\)は「スピードを少し変えたときに、エネルギーの無駄がどう変化するか」を見ています。

大自然は、「位置を変えたい要求」と「スピードを変えたい要求」のバランス（差）が常に「ゼロ」になるような、最も無駄のない絶妙な妥協点を探しながら物体を動かしているのです。

そして数学的に最も重要なのは、左側の項の先頭に\(\frac{d}{dt}\)（時間でもう一度微分する） という記号がついていることです。「速度（\(v\)）」という1階微分されたものを、さらに時間で微分するとどうなるでしょうか？ 答えは\textbf{「加速度（2階微分）」になります。 つまり、オイラー＝ラグランジュ方程式は、エネルギーのバランスを計算しているように見えて、結果的にはニュートンの\(F=ma\)と同じ「加速度（2階微分）を求める方程式」}になるのです。エネルギーというスカラーから出発しても、最終的に弾き出される物体の動きはニュートン力学と完璧に一致します。

\section{同じ現象、全く違う世界観：デカルト座標からの解放と大局的最適化}

ボールを空に向かって投げたときの美しい放物線を計算する際、「ニュートンの運動方程式」を使っても、「最小作用の原理（オイラー＝ラグランジュ方程式）」を使っても、出てくる答え（軌道）は数学的に全く同じになります。

しかし、その背後にある「宇宙の捉え方（パラダイム）」は天地ほど異なります。それは空間と時間をどのように認識するかという、極めて哲学的な世界観の衝突でもあります。

\subsection{デカルト座標の呪縛と「一般化座標」による解放}

17世紀、ルネ・デカルトは空間を縦・横・高さ（\(x,y,z\)）の直交するマス目で切り刻み、代数（数式）で幾何学（図形）を表現する\textbf{「デカルト座標系」}を発明しました。ニュートン力学は、この真っ直ぐな方眼紙の上で力を矢印（ベクトル）として分解することで大成功を収めました。

しかし、これには弱点がありました。例えば、ひもに吊るされた「振り子」の動きを\(x,y\)座標で計算しようとすると、「ひもの長さが常に一定である（\(x^2+y^2=l^2\)）」という面倒な制約条件をいちいち数式に組み込まなければならず、計算が異様に複雑になってしまうのです。ニュートンのベクトル力学は、デカルトが引いた「直交する直線の目盛り」に強く縛られていました。

一方、エネルギー（スカラー量）をベースにする解析力学は、この\textbf{「座標系の呪縛」から物理学を見事に解放しました。オイラー＝ラグランジュ方程式は、デカルトのような直線距離の\(x,y\)座標だけでなく、振り子の「角度（\(\theta\)）」であっても、あるいはグニャグニャに曲がった曲面上の特殊な目盛りであっても、人間が計算しやすいように設定したどんな測り方（これを一般化座標}と呼びます）を用いても、全く同じ美しい数式の形を保ち続けるのです。

「宇宙の真理は、人間が勝手に引いた方眼紙の目盛り（座標）の取り方などには左右されない」。

解析力学がもたらしたこの\textbf{「座標変換に対する不変性」}という洗練された視点は、のちにアインシュタインが「曲がった時空」を記述する一般相対性理論を生み出す決定的な土台となりました。さらには、現代のAI（多様体学習：Manifold Learning）が、人間には理解不能なほど歪んだ高次元のデータ空間から、座標軸に依存しない「本質的な特徴」だけを抽出する際の強力な数学的基盤ともなっています。

\subsection{微分と積分：時間を「切り刻む」か「足し合わせる」か}

空間の捉え方だけでなく、時間の捉え方においても両者は劇的に対立します。

ニュートンの力学は\textbf{「微分」の力学です。微分とは、時間を極限まで細かく切り刻むことです。ニュートンの視点では、ボールは「今の瞬間の位置」と「今の瞬間に受けている重力」という足元の情報しか持っていません。その足元の情報（原因）によって、ほんの0.001秒先のわずかな移動（結果）が決まります。これを無限に繰り返すことで軌道が作られます。目隠しをして、一歩ずつ足元の傾きだけを頼りに進んでいく、「局所的（Local）」}な視点です。

一方、解析力学は\textbf{「積分」の力学です。積分とは、時間をスタートからゴールまで全体にわたって足し合わせることです。ラグランジュの視点では、ボールは途中の瞬間瞬間を気にしません。「スタート地点」から「ゴール地点」までの全体のシナリオを上空から俯瞰し、考えうるすべてのルートのエネルギーの合計（作用）を計算した上で、最も無駄のないルートを一発で見つけ出して進んでいく、「大域的（Global）」}な視点なのです。

\subsection{因果律 vs 目的論：「背中を押される」か「未来に導かれる」か}

この数学的なアプローチの違いは、時間の流れに対する劇的な解釈の違いを生み出します。

ニュートン力学は、過去が現在を決め、現在が未来を決めるという\textbf{「因果律（原因と結果）」}に支配されています。ボールは常に「過去からの力に背中を押されて」前に進みます。

しかし最小作用の原理は、まるでボールが「あそこのゴールに、この時間で到着する」という未来の目的を最初から知っており、それに合わせて現在の進み方を決定しているかのように見えます。これを\textbf{「目的論的」}な振る舞いと呼びます。ボールは過去から押されているのではなく、未来のゴールから「最も無駄のないルートで来なさい」と見えざる手で導かれているようにすら錯覚させられます。

ただし、これは「未来が現在を本当に支配している」という意味ではありません。ニュートンの運動方程式とオイラー＝ラグランジュ方程式は、同じ運動を別々の言葉で記述しているだけです。前者は「いま働く力」から一歩先を計算し、後者は「始点と終点を結ぶ経路全体」を条件として眺めます。因果律を壊しているのではなく、同じ現象を局所的にも大域的にも記述できるところに、解析力学の深さがあります。

\subsection{宇宙が隠し持つ「二面性」の奇跡と、さらなる深淵へ}

局所的には「一歩ずつ手探りで進む盲目の機械（因果律）」のように振る舞う大自然が、大域的には「無駄を極限まで嫌い、最も効率的なルートをあらかじめ知っている究極の怠け者（目的論）」として振る舞う。

この「全く異なる2つの世界観」が、数学的には寸分違わず完璧に一致するという事実に、当時の物理学者たちは戦慄し、そこに宇宙の創造主の途方もないエレガンス（数学的美しさ）を見出しました。

（なぜ、大自然はそのような「神の視点（全体最適）」を持っているように振る舞えるのでしょうか？ その恐るべき真の理由は、約200年後、20世紀の「量子力学（ファインマンの経路積分）」によって明かされることになります。「実はボールは、ただ一つの最適ルートを神のように見抜いているのではなく、『宇宙に存在する考えうるすべてのルートを同時に分身して通り抜け、その結果として最適ルートだけが残る』」という、さらに想像を絶するミクロの真理が隠されていたのです。これについては第7章以降の量子力学で詳しく触れます。）

このように、同じ物理現象であっても「エネルギーの大域的な最適化」と「座標系に依存しない空間の視点」を持つことで、物理学は全く新しい次元へと跳躍したのです。

\section{ハミルトンの大手術：総エネルギーと「位相空間」の魔法}

ラグランジュが完成させた解析力学という強力なツールがすでにあったにもかかわらず、19世紀に入り、アイルランドの天才数学者ウィリアム・ローワン・ハミルトンは、全く新しい\textbf{「ハミルトニアン（\(H=K+U\)）」}という概念を導入しました。

なぜ、わざわざそんなことをする必要があったのでしょうか？ そこには、物理学と数学をさらに上の次元へと引き上げる「3つの極めて実用的、かつ決定的な理由」がありました。このハミルトンの飛躍こそが、のちの量子力学への扉を開き、同時に現代のAI（人工知能）が世界を認識するための「高次元空間の数学」の直接的なルーツとなっていきます。

\subsection{理由1：複雑な「2階微分」を、シンプルな「1階微分」に分解する}

先ほど見たオイラー＝ラグランジュ方程式には、一つだけ数学的に厄介な特徴がありました。それは数式の中に「加速度（速度をさらに時間で微分した2階微分）」が含まれてしまうことです。2階微分方程式は、ボールの軌道が複雑になると数学的に解くのが極めて困難になります。

そこでハミルトンは、「位置（\(q\)）」と「運動量（\(p\)）」という2つの変数を使って、この厄介な1つの2階微分方程式を、「2つのシンプルな1階微分方程式」にパカッと分解してしまったのです。これを\textbf{「ハミルトンの正準方程式」}と呼びます。

\[
\frac{dq}{dt}=\frac{\partial H}{\partial p},\quad \frac{dp}{dt}=-\frac{\partial H}{\partial q}
\]

（※\(q\)は位置、\(p\)は運動量、\(H\)はハミルトニアン）

左側の式は「位置の変化スピード（速度）」、右側の式は「運動量の変化スピード（力）」を表しています。どちらも\textbf{1階微分（\(\frac{d}{dt}\)）}でピタリと止まっており、2階微分（加速度）という厄介者が綺麗に消え去っています。これにより、数学的な見通しと計算のしやすさが劇的に向上しました。

\subsection{理由2：位置と運動量を平等に扱う究極の抽象化（位相空間の誕生）}

ニュートンやラグランジュの世界では、主役はあくまで「位置（どこにいるか）」であり、「運動量（スピード）」はそのオマケ（位置が時間でどう変わるか）に過ぎませんでした。

しかし上記のハミルトン方程式を見ると、位置\(q\)と運動量\(p\)が、まるで鏡合わせのように全く対等で独立した主人公として扱われています。

ハミルトンは、この位置と運動量をそれぞれ縦軸と横軸にとった、現実の3次元空間とは異なる、抽象的な多次元空間を発明したのです。これを\textbf{「位相空間（Phase Space：相空間とも呼ばれる）」と呼びます。 この位相空間という新しいキャンバスの上では、どんなに複雑に動き回る振り子や惑星の動きも、「たった1つの点が、ハミルトニアン（総エネルギー）の等高線に沿って、滑らかな川の水のように流れていく幾何学的な軌跡」}として極めてシンプルに描くことができます。

重要なのは、モーペルテュイが夢見た\textbf{「変分原理（最小作用の原理）」は、ラグランジュの現実空間のルート探しだけでなく、このハミルトンの抽象的な「位相空間上のルート探し」においても全く同じように使える（ハミルトンの原理）}ということです。大自然は、どんな高次元の抽象空間であっても「無駄を最小化する最適化のルール」を貫いているのです。

\subsection{理由3：のちの「量子力学」を作るための完璧な設計図}

歴史的な結果論でもありますが、これがハミルトンの最大の功績です。

20世紀に入り、シュレディンガーやハイゼンベルクといった天才たちがミクロの波の法則（量子力学）を作ろうとしたとき、ニュートン力学やラグランジュ力学の数式をそのまま変換してもうまくいきませんでした。

しかし、ハミルトンが作った「ハミルトニアン（総エネルギー）」と「位置と運動量の対等な関係」の数式だけは、そのまま変数を『オペレーター（微分演算子）』に置き換えるだけで、見事に量子力学の数式（シュレディンガー方程式\(\hat{H}\psi=E\psi\)）へとピタリとはまったのです（これを正準量子化と呼びます）。ハミルトンの力学は、無意識のうちに次世代の物理学への完璧な架け橋を用意していました。

\section{AIの弱点と「全体を見渡す」ことの難しさ：高次元の迷宮}

\subsection{物理学からAIの高次元パラメータ空間へ}

さて、ここで現代のAI（人工知能）に話を戻しましょう。

この「抽象的な高次元の位相空間を用意し、その中でエネルギーが低い場所へと状態が流れていく様子を幾何学的に捉える」というハミルトンのアプローチは、約150年の時を経て、現代のディープラーニング（AI）を理解するための強力な比喩として蘇っています。

現代の大規模なAIは、数千万から数千億個という途方もない数のパラメータ（神経細胞の繋がりの強さ）を持っています。AIが学習によって少しずつ賢くなっていくプロセスは、物理学的な視点から見れば、まさにこの\textbf{「数千億次元という人間には想像もつかない超高次元の空間（位相空間のAI版）」}の中で、損失関数（AIの「間違いの量」を表す、エネルギー関数に似た量）の等高線を描き、エラーが低くなる谷底を目指して、現在の状態を表す「1つの点」を滑り落としていくプロセスとして理解できます。

ハミルトンが考案した「位相空間」という強力な数学的メガネがなければ、情報科学者もAIの高次元の学習プロセスが作り出す山や谷（損失地形：Loss Landscape）を数学的に取り扱うことはできなかったでしょう。大自然のエネルギーの振る舞いと、知能が学習するプロセスは、この「多次元空間における最適化」という共通のパラダイムで見事に結びついているのです。

\subsection{局所的アプローチの限界：「浅い水たまり」と「鞍点（馬の鞍）」の罠}

第1章で見たように、AIを賢くする心臓部である「勾配降下法」は、まさに\textbf{ニュートン力学的な「局所的アプローチ（微分の力学）」}です。AIは、数千億次元の複雑な地形に目隠しで降ろされ、足元の傾き（微分）だけを頼りに、一歩ずつ下り坂の方へと進んでいきます。

しかし、この目隠しによる局所的アプローチには、致命的な弱点があります。それが\textbf{「局所的最適解（Local Minima）に囚われる」}という問題です。

本当の一番深い谷底（大域的最適解：Global Minimum）は高い山の向こう側にあるのに、AIはたまたま足元にあった「浅い水たまり（窪地）」に落ち込むと、前後左右どこを向いても上り坂になっているため、「ここが世界で一番深い谷底に違いない」と勘違いして、そこで学習を止めてしまうのです。全体を見渡す俯瞰的な目を持たないAIは、「この山を一度登って越えれば、もっと素晴らしい正解がある」ということに決して気づけません。

さらに、数千億次元という超高次元空間ならではの、より恐ろしい罠が存在します。それが\textbf{「鞍点（Saddle Point）」}の問題です。

馬の背中に乗せる「鞍（くら）」を想像してください。ある方向（前後）から見るとそこは「谷底（極小）」に見えますが、別の方向（左右）から見るとそこは「山の頂上（極大）」になっています。

2次元や3次元の空間なら鞍点は珍しいですが、数千億次元の空間になると、確率的に「すべての方向が上り坂になっている本物の谷底（極小値）」よりも、「一部の方向だけ下り坂になっている鞍点」の方が圧倒的に多く発生します。AIが足元の傾きだけを頼りに進むと、この無数にある鞍点の平らな部分に迷い込み、どちらへ進めばいいか分からなくなって学習が完全に停滞してしまう（勾配消失）という悲劇が起こるのです。

\subsection{物理学の知恵を借りる：「モメンタム」による突破}

では、AIにも解析力学（変分法）のように「上空から全体を見渡して、一発で最適ルートを決める（大域的最適化）」アプローチを取らせれば良いのではないでしょうか？ 残念ながら、数千億個のパラメータが作り出す無数のルートをすべて計算して比較するには、宇宙の寿命よりも長い計算時間が必要になり、現在のコンピュータでは絶対に不可能です。

そこでAI研究者たちは、再び物理学から強力なアイデアを借りてきました。それが\textbf{「モメンタム（Momentum：運動量）」}の導入です。

ただの勾配降下法は「羽根のように軽いボール」が転がるようなもので、浅い水たまりや鞍点ですぐに止まってしまいます。しかし、ボールに「重さ（質量）」と「転がる勢い（運動量＝モメンタム）」を持たせたらどうでしょうか？ ボールは勢い余って小さな窪地をゴロゴロと飛び越え、鞍点もそのままのスピードで駆け抜け、より深く広い「本当の谷底」へと到達しやすくなります。

現在、世界中のディープラーニングで標準的に使われている最適化アルゴリズム（Adamなど）は、まさにこの「物理的な運動量や慣性の法則」を数式に組み込むことで、高次元空間の迷宮を力強く突破しているのです。

\section{強化学習と「未来の報酬」：物理とAIの歴史的合流}

空間的な「全体（大域的最適化）」を見渡すことが難しいなら、せめて時間的な「全体（未来）」を見渡して最適な行動を決定できないか？

この長期的視点をAIに持たせようとする試みが、\textbf{「強化学習（Reinforcement Learning）」}と呼ばれる分野で極めて大きな成果を上げています。ここには、物理学の解析力学と、AIの最適化理論が完全に合流する、数学史上のドラマチックな軌跡が隠されています。

\subsection{AIが「未来の報酬」を最大化する}

囲碁の世界チャンピオンを打ち破った「AlphaGo」や、ロボットの自律歩行などに使われる強化学習では、AIは「目の前のこの一手を打てば、今すぐ10点がもらえる」という局所的な（ニュートン的な）視点だけで行動しません。彼らは、ゲームが終了するまでに得られるであろう\textbf{「未来の報酬の合計（期待収益：Return）」を最大化（最適化）」}するように学習します。

時には、今すぐ手に入る10点をわざと見送り、しゃがんで相手にポイントを譲ることで、100手後に1000点という巨大なリターンを得るような「捨て石」を打つことすらあります。AIは目先の利益（足元の傾き）に惑わされず、ゲーム全体（スタートからゴールまで）のシナリオを見据えた「目的論的」な振る舞いを獲得するのです。

これは、ラグランジュが提唱した「途中は少しエネルギーが高くなっても、スタートからゴールまでのエネルギーの合計（作用積分）を最小化するルートを選ぶ」という解析力学の考え方と非常に深く共鳴しています。

\subsection{ベルマン方程式とハミルトン：歴史的・数学的な奇跡}

この強化学習の「未来の報酬の最適化」を支えている絶対的な数学の土台が、1950年代に数学者リチャード・ベルマンが考案した「動的計画法」と\textbf{「ベルマン方程式」}です。ベルマン方程式は、「未来のゴールから逆算して、今のステップで取るべき最善の行動を決定する」という魔法のような漸化式です。

ただし、強化学習のAIも未来を実際に見ているわけではありません。価値関数と呼ばれる「この状態にいると、将来どれくらい得をしそうか」という見積もりを更新し続けているだけです。つまり、ここでも目的論に見える振る舞いの正体は、未来全体を一つの量として評価する数学的な仕組みなのです。

実は、このベルマン方程式は、AIのために突然ゼロから生み出されたものではありません。

その源流を辿ると、19世紀にハミルトン（とヤコビ）が解析力学を完成させるために導き出した\textbf{「ハミルトン・ヤコビ方程式」}に行き着きます。ハミルトン・ヤコビ方程式は、最小作用の原理に従って「目的地までの残りコスト」を最小化する物理学の究極の方程式です。

20世紀の数学者たちは、この物理学の方程式を「ロケットの燃料を最小限にして月に到達するにはどうすればいいか」といった工学的な最適制御理論（ハミルトン・ヤコビ・ベルマン方程式：HJB方程式）へと応用・拡張しました。そしてそれが、現代のAIエージェントがゲームや現実世界で「未来の報酬を最大化する行動」を学習するための根幹理論（強化学習）へと、シームレスに受け継がれていったのです。

「目隠しをして足元だけを見るAI（ニュートン的な勾配降下法）」から、いかにして「大局的な地形を想像させ、長期的な目的論的最適化（解析力学的な強化学習）を行わせるか」。

物理学がベクトルからスカラーへ、局所から大域へとパラダイムシフトを果たしたように、現代のAI研究もまた、物理学が敷いた数学のレールの上を走りながら、その知性の視座を「全体最適」へと引き上げようと格闘し続けています。

次章では、この「エネルギーと最適化」の物語に、まったく新しい強烈な概念が乱入してきます。それが、19世紀の産業革命とともに生まれた「熱力学」、そして宇宙を支配する冷酷なる法則「エントロピー」です。

\section*{【第2章 章末ディスカッション課題】}

大自然が「最小作用の原理」に従って動くとき、それは本当に「未来のゴールを知っている（目的論）」のでしょうか？ それとも、ただの数学的な結果がそう見えているだけなのでしょうか？

今のAIは「足元の傾き（微小な変化）」だけを見て学習しています。もし、大自然のように「全体のシナリオを一瞬で計算し尽くす」ことができる究極のコンピュータ（例えば量子コンピュータなど）が完成したとしたら、AIの「知能の形」はどのように変化すると思いますか？


\section*{参考文献（学術誌）}
\begin{enumerate}[leftmargin=2.4em]
\item Gray, C. G. and Taylor, E. F. ``When action is not least.'' \textit{American Journal of Physics}, 75(5), 434--458, 2007. DOI: \url{https://doi.org/10.1119/1.2710480}.
\item Papastavridis, J. G. ``The principle of least action as a Lagrange variational problem: Stationarity and extremality conditions.'' \textit{International Journal of Engineering Science}, 24(8), 1437--1443, 1986. DOI: \url{https://doi.org/10.1016/0020-7225(86)90072-8}.
\item Sutton, R. S. ``Learning to predict by the methods of temporal differences.'' \textit{Machine Learning}, 3, 9--44, 1988. DOI: \url{https://doi.org/10.1007/BF00115009}.
\item Mnih, V. et al. ``Human-level control through deep reinforcement learning.'' \textit{Nature}, 518, 529--533, 2015. DOI: \url{https://doi.org/10.1038/nature14236}.
\end{enumerate}


